\section{The classical domain inequality and sharpness}
\label{sec:classical}

We prove Corollary~\ref{cor:cartan-hadamard} and verify the optimality
of the filling coefficient. Both conclusions use the absence of
compactly supported top-dimensional cycles in a connected noncompact
oriented manifold.

\begin{lemma}\label{lem:top-dimensional-constancy}
Let $M$ be a connected, noncompact, oriented Riemannian manifold of
dimension $d$. If $C\in\I_d(M)$ has compact support and
$\partial C=0$, then $C=0$.
\end{lemma}

\begin{proof}
In a relatively compact oriented coordinate ball, the
top-dimensional representation of an integral current is integration
against an integer-valued locally integrable coefficient $\theta$.
Testing $\partial C=0$ with compactly supported smooth
$(d-1)$-forms shows that every distributional partial derivative of
$\theta$ is zero. Thus $\theta$ is constant almost everywhere on that
ball. For completeness, convolution with a smooth kernel gives
functions with zero ordinary gradient on each smaller ball; passing
to the local $L^1$ limit proves the assertion. Compatibility on
overlaps and connectedness give one constant on $M$. Some nonempty
open set lies outside the compact support, so this constant is zero.
\end{proof}

\begin{proof}[Proof of Corollary~\ref{cor:cartan-hadamard}]
By the Cartan--Hadamard theorem and Hopf--Rinow, $M$ is a proper
CAT(0) space and is diffeomorphic to $\R^d$; in particular it is
oriented and noncompact. See~\cite[Sections~I.3, II.1A and II.4]{BridsonHaefliger1999}
for the metric comparison and completeness statements.
Let $[[\Omega]]$ be its integration
current. Its mass is $\vol(\Omega)$, and the mass of
$T=\partial[[\Omega]]$ is $\operatorname{area}(\partial\Omega)$.
Theorem~\ref{thm:main}, with $n=d-1$, gives a compactly supported
$S\in\I_d(M)$ with $\partial S=T$ and
\[
 \M(S)\le
 \frac{\operatorname{area}(\partial\Omega)^{d/(d-1)}}
 {d^{d/(d-1)}\omega_d^{1/(d-1)}}.
\]
Lemma~\ref{lem:top-dimensional-constancy} applies to
$S-[[\Omega]]$, so $S=[[\Omega]]$. Rearranging its mass bound
proves~\eqref{eq:classical}.

For a relatively compact set $\Omega$ of finite perimeter, the
standard representation of its integration current gives
$[[\Omega]]\in\I_d(M)$, with mass $\vol(\Omega)$ and boundary
mass equal to its perimeter; see~\cite[Chapter~3, Section~4, Theorem~4.3;
Chapter~6, Section~3, Remark~3.15]{Simon2014}. The same
argument applies verbatim.
\end{proof}

For $d=2$, the classical disk inequality gives
$L(\partial D)^2\ge4\pi\operatorname{area}(D)$ for a smooth
Riemannian disk of Gaussian curvature at most zero; see
\cite[Section~1.2, Theorem~3]{Izmestiev2014} for this formulation.
Filling the holes in a bounded smooth domain on a Cartan--Hadamard
surface increases area and decreases boundary length. Applying the disk
inequality to each resulting component, and then summing, proves the
same bound for the domain because $(\sum_j L_j)^2\ge\sum_jL_j^2$.
Together with Corollary~\ref{cor:cartan-hadamard}, this establishes the
Euclidean Cartan--Hadamard domain inequality in every ambient dimension.

To see that the coefficient in Theorem~\ref{thm:main} is optimal,
take $X=\R^{n+1}$ and let $T$ be the oriented boundary of the
Euclidean ball $B_R$. Its mass is $s_nR^n$, while
\[
 \M([[B_R]])=\omega_{n+1}R^{n+1}
 =c_n(s_nR^n)^{(n+1)/n}.
\]
Any compactly supported integral filling of $T$ equals $[[B_R]]$
by Lemma~\ref{lem:top-dimensional-constancy}. Hence equality occurs
and no smaller universal coefficient is possible.
